%=====================================================================
\section{Solving a test instance: \code{solve\_rl\_f4.py}}
\label{sec:solve}

\src{../code/solve_rl_f4.py}{81}{103}

\begin{enumerate}
  \item Line 84: one instance becomes a batch of size 1.
  \item Lines 85--89: $A=8$ map symmetries. The encoder runs once per symmetry (line 89): the 8 rotated or mirrored copies of the map give the network 8 different inputs for the same instance (distances do not change, so every plan is valid for all of them).
  \item Lines 86--91: for each symmetry, \code{per} $=1+64$ episodes: the first one \emph{greedy} (always the most likely action), the other 64 \emph{sampled}. In total $8\times65=520$ episodes.
  \item Lines 93--96: the episodes are played in chunks of up to 1\,040 at once (memory); every chunk is a call of \code{rollout}.
  \item Lines 97--102: the best episode is kept, by the key $10^{12}\times(\text{dropped cars})+Z$, i.e.\ fewest dropped cars first, then lowest $Z$; \code{plan\_of} turns it into a plan.
\end{enumerate}

\code{run\_one} (lines 111--150) then:
\begin{itemize}
  \item scores the plan with \code{results\_io\_f4.evaluate}, which calls the independent checker \code{check\_plan\_f4.validate} (every rule of the MILP, $Z$ recomputed from scratch, exactly as for MILP plans);
  \item compares the simulator's $Z$ with the checker's (line 130);
  \item writes one row in \code{runs.csv} (status \code{plan} or \code{invalid\_plan}, $Z_1$--$Z_6$, service numbers, solve time) and one row per car and per van in \code{plan.csv}.
\end{itemize}
Every instance has its own random seed (line 121), so results do not depend on the order of solving, and a resumed run gives the same plans. A failing instance becomes a row with status \code{error}; the run continues.

\inshort{Solving plays 520 episodes per instance (8 views of the map, 65 tries each), keeps the plan with the fewest unplaced cars and then the lowest cost, and scores it with the same checker as the MILP.}
